% ch-a5-02.tex —— 第 5 篇（Ben-Zvi–Chen–Helm–Nadler）续片，底本 PDF 446–457 = 印刷页 433–444

\ysrc{433}
\subsubsection{形式圈与幂幺圈}
前面我们看到，当 $X$ 是概形时，指数映射把圈空间与奇切丛等同起来。当 $X$ 是叠时，指数映射甚至没有良好定义：例如，不存在代数叠的映射 $\exp:\frakg/G\to G/G$。

为了写下指数映射，我们需要限制到\emph{形式圈空间}（formal loop space）$\wh\LL X$，即 $\LL X$ 沿常值圈 $X\subset\LL X$ 的完备化。由 [BN10a] 的定理 6.9，形式的奇切丛与形式圈空间之间有 $S^1$-等变等价：
\[
\exp:\ \wh{\mathbb T}_X[-1]\xrightarrow{\ \sim\ }\wh\LL X.
\]
再引入圈空间 $\LL X$ 的另一个中间变体较为方便，即\emph{幂幺圈空间}（unipotent loop space）：
\[
\LL^uX:=\Map(B\GG_a,X).
\]
注意 $B\GG_a$ 在自身上的作用赋予 $\LL^uX$ 一个典范的 $B\GG_a$-作用。仿射化映射 $S^1\to B\GG_a$ 诱导 $S^1$-等变映射 $\LL^uX\to\LL X$；在我们所处的设定中（$X$ 是带仿射对角的 Artin 1-叠），该映射是单态射（见 [Ch20] 的命题 2.1.24）。

总之，我们有一串 $S^1$-等变单态射：
\[
\wh\LL X\into\LL^uX\into\LL X.
\]

\begin{exam}{2.4}
我们沿用例 2.1 中的例子。

\begin{enumerate}[label=(\arabic*),leftmargin=2.6em,itemsep=0.2em]
\item 若 $X$ 是概形，则上述包含都是恒等，即
\[
\wh\LL X=\LL^uX=\LL X.
\]
也就是说，闭包含 $X\into\LL X$ 是导出的加厚。

\item 当 $X=BG$ 时，形式圈空间是 $\wh\LL(BG)=\wh G_e/G$，即形式群的伴随商；幂幺圈空间是 $\LL^u(BG)=\wh G_{\mathcal U}/G$，即沿幂幺锥作形式完备化所得的伴随商。

\item 映射 $Y/G\to BG$ 给出映射 $\LL(Y/G)\to G/G$。由 [Ch20] 的命题 2.1.20 与 2.1.25，形式圈空间是沿 $\{e\}/G$ 的逆像的完备化，幂幺圈空间是沿 $\mathcal U/G$ 的逆像的完备化。
\end{enumerate}
\end{exam}

我们用下图总结各种圈空间：
\begin{equation}
\tag{2.1}\label{eq:a5-02-1}
\begin{tikzcd}[cells={nodes={draw,inner sep=5pt}},row sep=3.2em,column sep=2.6em]
& \begin{tabular}{c}$X$ 是概形\\ $\wh\LL X=\LL^uX=\LL X$\end{tabular}\arrow[dl]\arrow[dr] & \\
\begin{tabular}{c}$X$ 是叠，带乘性\\ 自同构\\ $\wh\LL X=\LL^uX\subsetneq\LL X$\end{tabular}\arrow[dr] && \begin{tabular}{c}$X$ 是叠，带加性\\ 自同构\\ $\wh\LL X\subsetneq\LL^uX=\LL X$\end{tabular}\arrow[dl] \\
& \begin{tabular}{c}$X$ 是叠\\ $\wh\LL X\subsetneq\LL^uX\subsetneq\LL X$\end{tabular} &
\end{tikzcd}
\end{equation}

\ysrc{434}
\subsection{等变层的 Koszul 对偶描述}
我们想要研究的主要对象，是圈空间 $\LL X$ 上 $S^1$-等变凝聚层的范畴。所有技术性讨论都推迟到附录 A，这里只陈述将要引用的主要结果。特别地，关于我们将要打交道的范畴——（紧重正规化的）$S^1$-等变凝聚层 $\QCoh^!(\LL X)^{\omega S^1}$——的精确定义，见 §A.1。

\subsubsection{概形情形}
回忆 §2.2.2：当 $X$ 是概形时，Hochschild–Kostant–Rosenberg（或指数）映射给出等价
\[
\exp:\ \mathbb T_X[-1]\xrightarrow{\ \sim\ }\LL X,
\]
它与 $S^1$-作用经由 $B\GG_a$-作用的分解相容。于是有等价
\[
\QCoh^!(\LL X)^{\omega S^1}\xrightarrow{\ \sim\ }\QCoh^!(\mathbb T_X[-1])^{\omega B\GG_a}.
\]
还要回忆：右端的 $B\GG_a$-作用提升为 $B\GG_a\times\GG_m$-作用，因此我们同样可以考虑分次增强 $\QCoh^!(\mathbb T_X[-1])^{\omega B\GG_a\rtimes\GG_m}$。Kapranov [Ka91] 证明了下述版本的 Koszul 对偶，[BN10a] 的推论 5.2 用圈空间的语言对它作了解释。

\begin{thm}{2.5（概形圈空间的 Koszul 对偶）}\label{thm:a5-02-1}
设 $X$ 是光滑概形，$\DD(X)$ 是 $X$ 上 D-模的范畴，$F\DD(X)$ 是 $X$ 上滤过 D-模的范畴，即微分算子上阶滤过所附的 Rees 代数上的模。

此时我们有相容的等价：
\begin{center}
\begin{tikzcd}[row sep=3.2em,column sep=1.2em]
\QCoh^!(\mathbb T_X[-1])^{\mathrm{Tate}\rtimes\GG_m}
& \QCoh^!(\mathbb T_X[-1])^{\omega B\GG_a\rtimes\GG_m}\arrow[l,"\sim"']\arrow[r,"\sim"]\arrow[d,"\sim"]
& \QCoh^!(\mathbb T_X[-1])^{\GG_m}\arrow[d,"\sim"]\\
\DD(X) & F\DD(X)\arrow[l]\arrow[r] & \QCoh^!(\mathbb T_X^*)^{\GG_m}
\end{tikzcd}
\end{center}
\end{thm}

\subsubsection{叠情形}
当 $X$ 是叠时，仅靠 Koszul 对偶只足以描述形式圈空间 $\wh\LL X$ 或幂幺圈空间 $\LL^uX$ 上的 $S^1$-等变凝聚层。（回忆：当 $X$ 是概形时，包含 $\wh\LL X\simeq\LL^uX\simeq\LL X$ 都是等价。）下述 Koszul 对偶是 [Ch23] 的定理；所涉及的「重正规化」范畴见 §A.2.1。特别地，范畴
\[
\widehat{\QCoh}^!(\wh\LL X)=\Ind\bigl(\widehat{\Coh}(\wh\LL X)\bigr)
\]
是 ind-凝聚层的一个适当重正规化的版本。

\begin{thm}{2.6（叠圈空间的 Koszul 对偶）}\label{thm:a5-02-2}
设 $X$ 是光滑几何叠，$\check\DD(X)$ 是 $X$ 上适当重正规化的 D-模范畴，$F\check\DD(X)$ 是 $X$ 上适当重正规化的滤过 D-模范畴。

此时我们有相容的等价：
\begin{center}
\begin{tikzcd}[row sep=3.2em,column sep=1.2em]
\widehat{\QCoh}^!(\wh\LL X)^{\mathrm{Tate}\rtimes\GG_m}
& \widehat{\QCoh}^!(\wh\LL X)^{\omega B\GG_a\rtimes\GG_m}\arrow[l,"\sim"']\arrow[r,"\sim"]\arrow[d,"\sim"]
& \widehat{\QCoh}^!(\wh\LL X)^{\GG_m}\arrow[d,"\sim"]\\
\check\DD(X) & F\check\DD(X)\arrow[l]\arrow[r] & \QCoh^!(\mathbb T_X^*)^{\GG_m}
\end{tikzcd}
\end{center}
\end{thm}

\ysrc{435}
特别地，我们有 Koszul 对偶（§A.1.3 意义下）：
\[
\widehat{\QCoh}^!(\wh\LL X)^{\mathrm{Tate}}\underset{\mathrm{Kos}}{\sim}\check\DD(X)\otimes\QCoh(\GG_m).
\]
进一步，我们不必过渡到形式圈空间 $\wh\LL X$，而可以更直接地与幂幺圈空间 $\LL^uX$ 打交道。下述将发表于 [BCHN] 的定理说明了这一点；一个算例见例 A.14。

\begin{thm}{2.7}\label{thm:a5-02-3}
设 $X$ 是几何叠。拉回函子诱导等价
\[
\widehat{\QCoh}^!(\LL^uX)^{\mathrm{Tate}\rtimes\GG_m}\simeq\widehat{\QCoh}^!(\wh\LL X)^{\mathrm{Tate}\rtimes\GG_m},
\]
从而有等价
\[
\widehat{\QCoh}^!(\LL^uX)^{\mathrm{Tate}\rtimes\GG_m}\simeq\check\DD(X).
\]
\end{thm}

\subsection{Jordan 分解与等变局部化}
有了关于幂幺圈空间 $\LL^uX$ 上 $S^1$-等变凝聚层的前述理论，我们希望更进一步，描述整个圈空间 $\LL X$ 上的 $S^1$-等变凝聚层。我们的做法是把圈分解为半单部分与幂幺部分，然后对幂幺部分套用先前的理论，并用半单部分来扭。

\subsubsection{圈的 Jordan 分解}
本节专门讨论 $G$ 是约化群（更一般地，是中性分支为约化的仿射代数群）时的整体商叠 $X/G$ 情形。更详细的讨论见 [Ch20] §3。

考虑「特征多项式」映射
\[
\mu:\ \LL(BG)=G/G\lra G\sslash G,
\]
它从伴随商叠（adjoint-quotient stack）\tn{底本第 435 页此处作「adoint-quotient stack」，疑为「adjoint-quotient stack」之误。}到仿射商概形，即到参数化半单共轭类的簇。注意恒等元的类的原像恰是幂幺锥 $\mathcal U_G\subset G$。

对任意半单元 $\alpha\in G$，置 $[\alpha]:=\mu(\alpha)\in G\sslash G$，并以 $\mathbb O_\alpha=G\cdot\{\alpha\}\simeq G/G(\alpha)\subset G$ 记 $\alpha$ 的共轭类，其中 $G(\alpha)\subset G$ 是 $\alpha$ 的中心化子。于是我们得到群元素（视为圈）的 Jordan 分解
\[
\mu^{-1}([\alpha])\simeq G\times^{G(\alpha)}\{su\in G\mid s\in\mathbb O_\alpha,\ u\in\mathcal U_G,\ su=us\}.
\]
现在设概形 $X$ 带有 $G$-作用。取自然映射
\[
\LL(X/G)\lra\LL(pt/G)\simeq G/G\lra G\sslash G
\]
的纤维，就可以定义圈空间的各种变体。特别地，在 $G\sslash G$ 的点上取完备化，给出幂幺圈空间的推广；在 $G/G$ 的闭轨道上取完备化，给出形式圈空间的推广。

\begin{defn}{2.8}\label{defn:a5-02-1}
设 $G$ 是线性代数群，其中性分支约化，$G$ 作用在预叠 $X$ 上。以 $\mathbb O\subset G$ 记一个半单伴随轨道，它定义了闭子叠 $\mathbb O/G\subset G/G\simeq\LL(BG)$；以 $[\alpha]=:\mu(\mathbb O)\in G\sslash G$ 记它在仿射商中的类。
\end{defn}

\ysrc{436}
我们定义如下圈空间：

\begin{enumerate}[label=(\arabic*),leftmargin=2.6em,itemsep=0.2em]
\item $\alpha$-幂幺圈空间 $\LL^u_\alpha(X/G)$ 是 $\LL(X/G)$ 沿 $[\alpha]\in G\sslash G$ 的原像的完备化，等价地，是沿轨道 $\mathbb O/G$ 的饱和化\footnote{即 $G/G$ 中包含 $\mathbb O$ 作为其唯一闭轨道的极大闭子叠。}的完备化。
\item $\alpha$-形式圈空间 $\wh\LL_\alpha(X/G)$ 是 $\LL(X/G)$ 沿轨道 $\mathbb O/G$ 的原像的完备化。
\item $\alpha$-特化圈空间 $\LL'_\alpha(X/G)=\LL(X/G)\times_{\LL(BG)}\mathbb O/G$ 是 $\LL(X/G)$ 在 $\mathbb O/G$ 上的（导出）纤维。
\end{enumerate}

尽管记号看起来并非如此，上述构造只依赖于轨道 $\mathbb O$（等价地，只依赖于 $[\alpha]\in G\sslash G$），而不依赖于代表元（或提升）$\alpha\in\mathbb O$ 的选取。一旦作这样的选取，我们就得到等价 $\mathbb O/G\simeq BG(\alpha)$，即 $\alpha$ 的中心化子的分类叠，于是可以定义下述（依赖于 $\alpha$ 的）圈空间：

\begin{enumerate}[label=(\arabic*),leftmargin=2.6em,start=4,itemsep=0.2em]
\item $\alpha$-扭圈空间\footnote{也称为由 $\alpha$ 在 $X$ 上的作用所定义的自映射的导出不动点。特别地，$X$ 上的 $\alpha$-扭圈空间只以该自映射为输入，而不需要整个 $G$ 在 $X$ 上的作用。} $\LL_\alpha(X)$ 是 $\LL(X/G)$ 沿 $\Spec k=\{\alpha\}\to BG(\alpha)\simeq\mathbb O/G\subset G/G$ 的（导出）纤维。
\end{enumerate}

$\LL(X/G)$ 上的 $S^1$-作用限制为 $S^1$-等变包含
\[
\LL'_\alpha(X/G)\subset\wh\LL_\alpha(X/G)\subset\LL^u_\alpha(X/G)\subset\LL(X/G).
\]
扭圈空间 $\LL_\alpha(X)$ 不容许 $S^1$-作用，但存在映射 $\LL_\alpha(X)\to\LL'_\alpha(X/G)$。

\subsubsection{圈的层与等变局部化}
这里陈述 [Ch20] 的等变局部化定理，它在给定半单部分 $\alpha$ 的情形下，用商叠 $X(\alpha)/G(\alpha)$ 上的幂幺圈来刻画商叠 $X/G$ 中的圈。这里 $X(\alpha)$ 表示关于由 $\alpha$ 生成的闭子群 $A=\langle\alpha\rangle\subset G$ 的同伦不动点 $X^A$。更详细的讨论见 §A.2.2。

考虑自然映射
\[
\ell_\alpha:\ \LL_\alpha(X(\alpha)/G(\alpha))\lra\LL_\alpha(X/G),
\]
它由对自然映射 $X(\alpha)/G(\alpha)\into X/G(\alpha)\to X/G$ 施用 $\LL_\alpha$ 得到。

\begin{thm}{2.9（导出圈空间的等变局部化）}\label{thm:a5-02-4}
设 $G$ 是线性代数群，其中性分支约化，$G$ 作用在局部诺特导出概形 $X$ 上，$\alpha\in G$ 是半单元。则幂幺 $\alpha$-局部化映射
\[
\ell^u_\alpha:\ \LL^u_\alpha(X(\alpha)/G(\alpha))\lra\LL^u_\alpha(X/G)
\]
是 $S^1$-等变的等价。对形式圈与特化圈上相应的局部化映射
\[
\hat\ell_\alpha:\ \wh\LL_\alpha\bigl(X(\alpha)/G(\alpha)\bigr)\lra\wh\LL_\alpha(X/G),\qquad
\ell'_\alpha:\ \LL'_\alpha\bigl(X(\alpha)/G(\alpha)\bigr)\lra\LL'_\alpha(X/G)
\]
同样成立。
\end{thm}

\begin{proof}
当 $X$ 光滑时，这就是 [Ch20] 的定理 A。一般情形由下述论证推出：按 loc. cit. 中的归约，可以假设 $\alpha$ 是中心元，并令 $A=\langle\alpha\rangle$。我们首先断言：每个导出概形都可以局部地写成光滑仿射 $A$-概形的迭代纤维积\tn{底本第 436 页此处作「interated fiber product」，疑为「iterated fiber product」之误。}。

\ysrc{437}
由于 $A$ 的中性分支是乘法环面，每个 $A$-概形都有由 $A$-闭开集组成的 Zariski 覆盖，因此可以假设 $X$ 是仿射的。又因 $X$ 是局部诺特的，可以选取 $\OO(X)$ 的一个 $A$-等变半自由消解，即：分次有限维 $A$-子表示 $V^*_0\subset\OO(X)$，它作为 dg 代数生成 $\OO(X)$；同样地，$V^*_1\subset\OO(V_0)$ 生成关系，$V_2\subset\OO(V_0\times V_1)$ 生成关系之间的关系\tn{底本第 437 页此处作 $V_2$，与前后文的 $V^*_1$、$V^*_0$ 体例不一，疑脱上标 $*$；照录原文。}，如此等等。这些拼成一个光滑仿射 $A$-概形的图表，其中竖直箭头是零包含，水平箭头由半自由消解中的微分给出：
\begin{center}
\begin{tikzcd}[row sep=2.6em,column sep=3.4em]
& V_0\arrow[d] & V_0\times V_1\arrow[d] & \cdots\arrow[d]\\
V_0\arrow[r] & V_0\times V_1\arrow[r] & V_0\times V_1\times V_2\arrow[r] & \cdots
\end{tikzcd}
\end{center}
该图表的极限是 $X$；而圈空间与同伦不动点都与极限交换，故可由光滑情形推出定理。
\end{proof}

最后，把等变局部化与 Koszul 对偶结合起来，得到下述结果。与上面一样，各种重正规化范畴的细节见 §A.2.1，扭 $S^1$-作用与 $\alpha$-平凡块的概念见 §A.2.3 与 §A.2.4。

\begin{thm}{2.10}\label{thm:a5-02-5}
设 $\alpha\in G$ 是半单元，$X/G$ 是整体商叠。我们有相容的等价
\begin{center}
\begin{tikzcd}[row sep=3.2em,column sep=1.1em]
\QCoh^!(\LL_\alpha(X/G))^{\mathrm{Tate}\rtimes\GG_m}_\alpha
& \QCoh^!(\LL_\alpha(X/G))^{\omega B\GG_a\rtimes\GG_m}_\alpha\arrow[l,"\sim"']\arrow[r,"\sim"]\arrow[d,"\sim"]
& \QCoh^!(\LL_\alpha(X/G))^{\GG_m}_\alpha\arrow[d,"\sim"]\\
\check\DD(X(\alpha)/G(\alpha))_\alpha
& F\check\DD(X(\alpha)/G(\alpha))_\alpha\arrow[l]\arrow[r]
& \QCoh\bigl(\mathbb T^*_{X(\alpha)/G(\alpha)}\bigr)_\alpha
\end{tikzcd}
\end{center}
它们对光滑拉回与紧合推前都是函子性的，分次版本亦然。

特别地，经由定理 2.7，我们得到 Koszul 对偶：
\begin{equation}
\tag{2.2}\label{eq:a5-02-2}
\widehat{\QCoh}^!\bigl(\LL^u_\alpha(X/G)\bigr)^{\mathrm{Tate}}_\alpha
\simeq\widehat{\QCoh}^!\bigl(\wh\LL(X/G)\bigr)^{\mathrm{Tate}}_\alpha
\underset{\mathrm{Kos}}{\sim}\check\DD\bigl(X(\alpha)/G(\alpha)\bigr)_\alpha.
\end{equation}
\end{thm}

\subsubsection{扭 Springer 层与分次提升}
现在考虑下述情形，它是我们迄今发展出的理论的一个例子。设 $\mu:\wt X/G\to X/G$ 是光滑叠的紧合映射，并令
\[
\mathcal S=\LL\mu_*\OO_{\LL(\wt X/G)}\in\Coh\bigl(\LL(X/G)\bigr).
\]
在 §3 中，这将是\emph{凝聚 Springer 层}。仍以 $\hat\ell_\alpha$ 记局部化映射，我们令
\[
\mathcal S(\alpha):=\hat\ell^!_\alpha\mathcal S\in\widehat{\Coh}\bigl(\wh\LL_\alpha(X(\alpha)/G(\alpha))\bigr)\subset\QCoh^!\bigl(\wh\LL_\alpha(X(\alpha)/G(\alpha))\bigr)
\]
为它到 $\alpha$-形式圈的 !-限制，它带有

\begin{enumerate}[label=(\arabic*),leftmargin=2.6em,itemsep=0.2em]
\item 一个典范的 $S^1$-等变结构，来自对偶化层 $\OO_{\LL(\wt X/G)}\simeq\omega_{\LL(\wt X/G)}$ 上典范的 $S^1$-等变结构（见 §A.1）；
\item 一个典范的 $\alpha$-平凡化：由于 $\alpha$ 在结构层上平凡地作用，$S^1$-等变结构可经由 $B\GG_a$-等变结构分解（见 §A.2.4）；
\end{enumerate}

\ysrc{438}
\begin{enumerate}[label=(\arabic*),leftmargin=2.6em,start=3,itemsep=0.2em]
\item 一个典范的分次提升 $\wt{\mathbf S}(\alpha)$，来自 $\omega_{\wh\LL(\wt X(\alpha)/G(\alpha))}$ 上的分次提升；后者来自把 $\wh\LL_\alpha(\wt X(\alpha)/G(\alpha))$ 与奇切丛等同起来（见 §A.1.3）。\footnote{尽管这些分次提升是典范的，它们却不能整体地作出，只能逐个参数地作出。}
\end{enumerate}

与 Koszul 对偶对应地，考虑诱导映射 $\mu^\alpha:\wt X(\alpha)/G(\alpha)\to X(\alpha)/G(\alpha)$。对应于 $S^1$-等变层 $\wt{\mathbf S}(\alpha)$ 的对象是滤过 D-模
\[
\wt{\mathbf S}(\alpha)\simeq\mu^\alpha_*\OO_{\wt X(\alpha)/G(\alpha)},
\]
其中 $\OO_{\wt X(\alpha)/G(\alpha)}$ 带有典范滤过。把相应的无滤过 D-模记为 $\mathbf S(\alpha)$。这些对象有（导出）自同态代数
\[
\mathcal H(\alpha)=\End_{\wh\LL_\alpha(X/G)}\bigl(\mathcal S(\alpha)\bigr)^{S^1},\qquad
\mathbf H(\alpha)=\End_{X(\alpha)/G(\alpha)}\bigl(\mathbf S(\alpha)\bigr),
\]
以及相应的分次提升 $\wt{\mathcal H}(\alpha)$、$\wt{\mathbf H}(\alpha)$。由 Koszul 对偶等价立刻得到下述推论。

\begin{cor}{2.11}\label{cor:a5-02-1}
把层 $\mathcal S(\alpha)$、$\wt{\mathbf S}(\alpha)$ 视为 $S^1$-等变范畴中的对象。我们有分次代数的等价 $\wt{\mathcal H}(\alpha)^\natural\simeq\wt{\mathbf H}(\alpha)$，以及 $\GG_m$-等变范畴的等价
\[
\langle\wt{\mathbf S}(\alpha)\rangle^{\GG_m,\natural}\simeq\Mod^{\GG_m}\bigl(\wt{\mathcal H}(\alpha)\bigr)^\natural\simeq\Mod^{\GG_m}\bigl(\wt{\mathbf H}(\alpha)\bigr)\simeq\langle\wt{\mathbf S}(\alpha)\rangle^{\GG_m},
\]
从而有 Koszul 对偶：
\[
\langle\mathcal S(\alpha)\rangle\simeq\Mod\bigl(\mathcal H(\alpha)\bigr)\underset{\mathrm{Kos}}{\sim}\Mod\bigl(\mathbf H(\alpha)\bigr)\simeq\langle\mathbf S(\alpha)\rangle.
\]
\end{cor}

\section{非阿基米德局部 Langlands 与圆作用}
\subsection{幂幺表示}
固定非阿基米德局部域 $F$，其剩余域为 $\F_q$，并固定 $F$ 上连通的、分裂的约化群 $G$（不过这里的许多讨论对非分歧群也适用）。Lusztig [Lu95b] 与 Solleveld [Sol18], [Sol19] 对这类群及其纯内型的幂幺表示建立了 Langlands 对应。（回忆：$G$ 的内型 $G^*$ 的表示 $\pi$ 称为幂幺的，是指存在 $G$ 的 parahoric 子群 $P$，以及 $P$ 的约化商的 $\F_q$-点上的幂幺表示 $\tau$，使得 $\pi$ 到 $P$ 的限制包含表示 $\tau$。）来自给定 $G^*$、$P$ 与 $\pi$ 的不可约幂幺表示，与 Hecke 代数 $\mathcal H_{G^*,P,\pi}$ 的不可约表示一一对应；该代数即所谓的「带不等参数的仿射 Hecke 代数」。

Lusztig 的构造给出如下双射：把数对 $(G^*,\pi)$（其中 $G^*$ 是 $G$ 的纯内型，$\pi$ 是不可约幂幺表示）与 $G$ 的\emph{幂幺扩张 Langlands 参数}的集合对应起来，后者即三元组 $(\sigma,n,\chi)$ 的 $\check G$-共轭类的集合，其中 $\sigma$ 是 $\check G$ 中的半单元，$n\in\NN$ 满足 $\sigma\cdot n=qn$，而 $\chi$ 是 $(\sigma,n)$ 的 $\check G$-中心化子 $\check G^{\sigma,n}$ 的分量群的一个不可约表示。

固定 $\check G$ 中的半单元 $\sigma$，就可以把 Vogan 簇 $\NN^{(\sigma,q)}/\check G^\sigma$ 定义为由满足 $\sigma\cdot n=qn$ 的 $n\in\NN$ 构成的空间；于是三元组 $(\sigma,n,\chi)$ 确定了 $\NN^{(\sigma,q)}/\check G^\sigma$ 中 $n$ 的轨道上的一个局部系统，从而是 $\NN^{(\sigma,q)}/\check G^\sigma$ 上的一个反常层。

\ysrc{439}
Lusztig 的论证把这类反常层与 Hecke 代数的表示联系起来：它把每个 Hecke 代数 $\mathcal H_{G^*,P,\pi}$ 在其中心各极大理想处的完备化，等同于 Vogan 簇 $\NN^{\sigma,q}/\check G^\sigma$ 上某些反常层的导出自同态。

我们把这一构造在 $G$ 的主块情形下具体化；所谓主块，即由 Iwahori 不动向量生成的 $G$ 的表示构成的满子范畴。这样的不可约表示恰是仿射 Hecke 代数 $\mathcal H_q:=\mathcal H_{G,B,\mathrm{triv}}$ 的不可约表示。该 Hecke 代数中心的极大理想与 $\check G$ 中的半单共轭类一一对应。

假设 $q$ 不是单位根，Lusztig 的构造把仿射 Hecke 代数 $\mathcal H_q(\sigma)$ 在对应于 $\sigma\in\check G$ 的极大理想处的完备化，实现为 $\NN^{(\sigma,q)}/\check G^\sigma$ 上某个可构造层的导出自同态的 Koszul 等价物。\footnote{特别地，前者处于上同调次数 0，后者在上同调次数 2 有生成元。讨论见 §3.4 与定义 A.10。}具体地，我们可以把 $\sigma$-Springer 层 $\mathbf S_q(\sigma)$ 定义为沿「$\sigma$-固定」Springer 消解
\[
\mu^{(\sigma,q)}:\ \wt\NN^{(\sigma,q)}/\check G^\sigma\lra\NN^{(\sigma,q)}/\check G^\sigma
\]
的推前 $\mu^{(\sigma,q)}_*\mathbb C_{\wt\NN^{(\sigma,q)}}$。

BBD 分解定理蕴含直和分解
\[
\mathbf S_q(\sigma):=\bigoplus_{(\mathcal O,\mathcal L)}E_{\mathcal O,\mathcal L}\otimes\mathbf{IC}(\mathcal O,\mathcal L)[d_{\mathcal O,\mathcal L}].
\]
出现在 $\sigma$-Springer 层 $\mathbf S_q(\sigma)$ 中的相交上同调层（由这样的局部系统延拓而来）正是给定 L-包中出现的主系列表示。此外，$\End(\mathbf S_q(\sigma))$ 是一个代数，其 0 次部分是仿射 Hecke 代数在由 $[\sigma]\in\check G\sslash\check G$ 给出的中心特征处的特化 $\mathcal H_{q,[\sigma]}$，因此 $\mathcal H_{q,[\sigma]}$ 作用在 $\mathbf S_q(\sigma)$ 的单直和项上。Deligne–Langlands 对应由这一分解给出，即 $E_{\mathcal O,\mathcal L}$ 是对应于由 $(\mathcal O,\mathcal L)$ 给出的 Langlands 参数的 $G(F)$-表示的 Iwahori 不动向量。

\subsection{幂幺 Langlands 参数的叠}
关于非分歧主系列的「代数」或「族」视角是在 [BCHN22] 中发展起来的。具体地，我们考虑\emph{幂幺 Langlands 参数}的叠，它可以用如下等价的方式定义。固定 $q$ 为素数的幂，从而是一个非零整数，精神上它是 $F$ 的剩余域的大小。

\begin{enumerate}[label=(\arabic*),leftmargin=2.6em,itemsep=0.2em]
\item 显式地，它是叠
\[
\mathbb L^u_{F,G}=\mathbb L^u_{q,G}:=\{n\in\NN_{\check G},\ g\in\check G\mid \mathrm{Ad}(g)\cdot n=qn\}/\check G.
\]
注意我们已经删去了 Deligne–Langlands 参数上「Frobenius 的像 $g$ 半单」这一条件。

\item 它是与 $q\in\GG_m$ 相联系的扭圈空间\footnote{第一个等同用到了 $q$ 不是单位根这一事实。}
\[
\LL_q\bigl(\wh\NN_{\check G}/\check G\bigr)=\LL_q\bigl(\check{\frakg}/\check G\bigr)
=\LL\bigl(\check{\frakg}/(\check G\times\GG_m)\bigr)\times_{\LL(B\GG_m)}\{q\}.
\]
\end{enumerate}

\ysrc{440}
由定义 2.8 的讨论，$q$-扭圈空间上没有 $S^1$-旋转。不过，存在来自 $\wh\NN_{\check G}/\check G=\LL^u(B\check G)$ 上作用的 $S^1$-作用，我们不考虑它。等价地，可以把该空间用 $\check G$ 中的幂幺元素（而非 Lie 代数中的幂零元）写出，此时 $q$ 通过指数作用。

\begin{enumerate}[label=(\arabic*),leftmargin=2.6em,start=3,itemsep=0.2em]
\item 它是 $q$-扭拓扑环面 $T_q$ 上 $\check G$-局部系统的叠 $\Loc_{\check G}(T_q)$ 的子叠\tn{底本第 440 页此处作「the stack of of $\check G$-local systems」，衍一 of。}，其中要求绕非扭圈的单项为幂幺；这里 $T_q$ 是把柱面 $S^1\times[0,1]$ 的两个边界沿次数 $q$ 的映射粘合所得的环面。空间 $T_q$ 在经线方向有 $S^1$-作用（通过「加速圈」），对应驯顺单项（我们不考虑它），但在纵向（Frobenius）方向没有。等价地，可以把 $T^u_q$ 定义为把 $B\GG_a\times[0,1]$ 沿「乘以 $q$」映射粘合所得的空间，于是 $\mathbb L^u_{q,G}=\Loc_{\check G}(T^u_q)$（这里 $q$ 不必是整数）。
\end{enumerate}

存在自然映射 $\mu:\mathbb L^u_{q,B}\to\mathbb L^u_{q,G}$，我们定义\emph{凝聚 Springer 层}为
\[
\mathcal S_q:=\mu_*\OO_{\mathbb L^u_{q,B}}\in\Coh\bigl(\mathbb L^u_{q,G}\bigr).
\]
我们有 [BCHN22] 的下述主要定理。

\begin{thm}{3.1}\label{thm:a5-02-6}
仿射 Hecke 代数自然同构于凝聚 Springer 层 $\mathcal S_q$ 的 Ext 代数，后者在非零上同调次数上消没。因此有全嵌入
\[
\mathcal H_q\text{-}\mathrm{mod}\xrightarrow{\ -\otimes_{\mathcal H_q}\mathcal S_q\ }\langle\mathcal S_q\rangle\xhookrightarrow{\ }\QCoh^{!}\bigl(\mathbb L^u_{q,G}\bigr),
\]
从 $G(F)$ 的非分歧主系列到幂幺 Langlands 参数叠上的 ind-凝聚层。
\end{thm}

因此，把我们对幂幺主系列表示的族化（或凝聚化、代数化）实现，与上述由 $\sigma$-Springer 理论在参数处给出的特化（或可构造化、拓扑化）实现联系起来，是一个自然的问题。

在 $G=\mathrm{GL}_n$ 的情形，Bushnell–Kutzko 的型理论结合 Harris–Taylor 与 Henniart 的局部 Langlands 定理，使我们能把 $G$ 的整个光滑表示理论归约到非分歧主系列。于是推出：$\mathrm{GL}_n$ 的整个光滑表示范畴全嵌入到所有 Langlands 参数的叠 $\mathbb L_{F,\mathrm{GL}_n}$ 上的 ind-凝聚层中，如 [H20, DHKM20] 所构造：

\begin{thm}{3.2（[BCHN22]）}\label{thm:a5-02-7}
存在全嵌入
\[
D(\mathrm{GL}_n(F))\xhookrightarrow{\ }\QCoh^{!}\bigl(\mathbb L_{F,\mathrm{GL}_n}\bigr),
\]
从 $\mathrm{GL}_n(F)$ 的光滑表示的导出范畴到 Langlands 参数叠上的 ind-凝聚层。该嵌入由下述性质刻画：把不可约尖点表示送到对应 Langlands 参数处的摩天层，并与抛物诱导相容。
\end{thm}

\ysrc{441}
\subsection{圈空间与循环形变}
幂幺 Langlands 参数叠 $\mathbb L^u_q$ 的一个有趣特征，是它与等变幂零锥的导出圈空间很接近：
\[
\LL(\check\NN/\check G\times\GG_m)\simeq\{n\in\check\NN,\ g\in\check G,\ q\in\GG_m:\ g\cdot n=qn\}/\check G\times\GG_m,
\]
而定理 3.1 有一个在整个圈空间上定义的变体，它把以 $q\in\GG_m$ 为参数的仿射 Hecke 代数恢复出来。更重要的是：叠 $\mathbb L^u_q$——即 $\{q\}\to\GG_m/\GG_m$ 上的纤维——并不继承圆作用，而\emph{分次幂幺 Langlands 参数}的叠
\[
\mathbb L^u_{\underline q}:=\mathbb L^u_q/\GG_m
\]
——即 $\{q\}/\GG_m\in\GG_m/\GG_m$ 上的纤维——则继承。此外，定理 3.1 在 $\mathbb L^u_{\underline q}$ 上原样成立，而且实际上有 $S^1$-等变的增强：在其中凝聚 Springer 块在等变参数 $k[u]=H^*(BS^1)$ 上平凡地形变。

\begin{thm}{3.3（[BCHN22]）}\label{thm:a5-02-8}
我们有全嵌入
\[
\mathcal H_q[u]\text{-}\mathrm{mod}\simeq\langle\mathcal S_q\rangle\subset\QCoh^{!}\bigl(\mathbb L^u_{\underline q}\bigr)^{\omega S^1},
\]
从 $G(K)$ 的非分歧主系列的 $k[u]$-基变换到分次幂幺 Langlands 参数叠上的（重正规化的）循环层。
\end{thm}

\subsection{与可构造层、分次 Hecke 代数的关系}
§2.1 的幂幺圈空间构造，把 $G$ 的表示论的这种统一的凝聚描述与固定无穷小特征处的反常描述联系起来。事实上，固定半单元 $\sigma\in\check G$，就可以考虑 $\mathbb L^u_q$ 沿下述配对构成的轨迹的形式完备化 $\mathbb L^u_{q,[\sigma]}$：这些配对 $(\sigma',n)$ 使 $\sigma'$ 的半单部分落在 $\sigma$ 的共轭类 $[\sigma]$ 中。然后可以把该形式完备化与 Vogan 簇在配对 $\alpha:=(\sigma,q)$ 处附带的 $q$-特化幂幺圈空间 $\LL^u_q(\check\NN^\alpha/\check G^\alpha)$ 等同起来。

于是定理 2.6 在幂幺 Langlands 参数叠上的（$S^1$-等变）凝聚层与 $\check\NN^\alpha/\check G^\alpha$ 上的 $\check G^\alpha$-等变\footnote{由于我们是在 $q\in\GG_m$ 处特化而非完备化，所得的 D-模是 $\GG_m$-可构造的而非 $\GG_m$-等变的。不过，在 $\check G^\alpha$-等变的前提下，$\GG_m$-可构造性是自动的；见 [Ch23] 的注 3.3.13。} D-模之间给出联系；该等价把凝聚 Springer 层到 $\mathbb L^u_{q,[\sigma]}$ 的 !-限制（以下记为 $\mathcal S(\alpha)$）送到对应于 $\alpha$-Springer 层 $\mathbf S(\alpha)$ 的 D-模。把推论 2.11 直接用于这一情形，得到下述结果。

\begin{thm}{3.4}\label{thm:a5-02-9}
设 $\mathbf H(\alpha)=\End_{\check\NN^\alpha/\check G^\alpha}(\mathbf S(\alpha))$，$\mathcal H(\alpha)$ 表示仿射 Hecke 代数在 $\alpha$ 处的完备化。我们有 Koszul 对偶：
\[
\widehat{\QCoh}^{!}\bigl(\mathbb L^u_{q,[\sigma]}\bigr)^{\omega S^1}\supset\langle\mathcal S(\alpha)\rangle
\simeq\Mod\bigl(\mathcal H(\alpha)\bigr)\underset{\mathrm{Kos}}{\sim}\Mod\bigl(\mathbf H(\alpha)\bigr)
\simeq\langle\mathbf S(\alpha)\rangle\subset\DD\bigl(\check\NN^\alpha/\check G^\alpha\bigr).
\]
按 2-周期的方式工作，我们得到满子范畴的等同
\[
\widehat{\QCoh}^{!}\bigl(\mathbb L^u_{q,[\sigma]}\bigr)^{\mathrm{Tate}}\supset
\Mod_{k[u,u^{-1}]}\bigl(\mathcal H(\alpha)^{\mathrm{per}}\bigr)\subset\DD\bigl(\check\NN^\alpha/\check G^\alpha\bigr)^{\mathrm{per}}.
\]
\end{thm}

\ysrc{442}
代数 $\mathbf H(\alpha)$ 是 Lusztig 分次 Hecke 代数 [Lu88, Lu89] 的完备化的剪移（shearing）。我们指出，Koszul 对偶中出现的关于分次提升的等变化、剪移与去等变化（如 §A.1.3 所讨论），与 [Sol18, Sol19] 中各陈述的相应分次相容；而且按 2-周期的方式工作可以避开这些复杂之处。

\subsection{周期循环层与纯内型}
\subsubsection{情形 $G=GL_n$}
在 $\mathrm{GL}_n$ 的情形，非分歧主系列穷尽了所有幂幺表示（且 $\mathrm{GL}_n$ 只有一个纯内型）。同样在 $\mathrm{GL}_n$ 的情形，Springer 理论大大简化：Springer 层生成幂零锥上所有等变的 D-模层，即出现在 Springer 层分解中的局部系统只有平凡的，而且所有这样的局部系统都出现。这给出 D-模范畴 $\DD(\NN_{\mathrm{GL}_n}/\mathrm{GL}_n)$ 与分次 Hecke 代数的模之间的导出等价。而且对任意中心特征同样成立，即对幂零锥的 $(\sigma,q)$-变体也成立。换言之，凝聚 Springer 层的像在任意参数处完备化后生成周期循环层的范畴。

在 [BCHN] 中我们证明如下结果：

\begin{thm}{3.5}\label{thm:a5-02-10}
对 $G=\mathrm{GL}_n$，嵌入
\[
D(\mathrm{GL}_n)\otimes_k k[u]\subset\QCoh^{!}\bigl(\mathbb L_{F,\mathrm{GL}_n}\bigr)^{\omega S^1}
\]
在把 $u$ 可逆化之后成为等价：
\[
D(\mathrm{GL}_n)\otimes_k k[u,u^{-1}]\simeq\QCoh^{!}\bigl(\mathbb L_{F,\mathrm{GL}_n}\bigr)^{\mathrm{Tate}}.
\]
\end{thm}

与定理 3.2 一样，标准技巧把断言归约到 Iwahori 块的相应陈述：在从所有凝聚层到周期循环层的形变下，只有凝聚 Springer 块存活。在那里我们证明，Springer 消解与 Hecke 代数具有足够强的有限性，使我们能从逐个完备化的相应生成断言推出 Springer 层的生成性
\[
\langle\mathcal S_q\rangle\simeq\QCoh^{!}\bigl(\mathbb L^u_{\underline q}\bigr)^{\mathrm{Tate}}
\]
（一种「断裂定理」）。

\subsubsection{一般 $G$ 的情形}
对一般的约化群，幂零锥及其 $(\sigma,q)$-不动轨迹上有更多的反常层，它们来自 $\check G$ 的 Levi 子群上的\emph{尖点局部系统}，Lusztig 作了分类（例如见 [Lu88, Lu95a]）。另一方面，$G$ 及其纯内型的幂幺表示也多于非分歧主系列。如 §3.1 所述，两者由 Lusztig [Lu95b] 对 $G$ 为伴随、单、非分歧情形所证明、并由 Solleveld [Sol18, Sol19] 推广到所有连通群的幂幺局部 Langlands 对应相匹配。

人们预期（沿 [Zhu25, He20, Fa16] 的诸猜想的变体，见 Zhu 即将发表的工作 [Zhu25]）：$\mathbb L^u_q$ 上凝聚层的整个范畴参数化了 $G$-等晶体的稳定化子族的幂幺表示。

\ysrc{443}
在这个大范畴内部，人们可以尝试找出\emph{纯子范畴}，它参数化 $G$ 自身的纯内型族的表示。对半单的 $G$，这些与基本等晶体（即稳定化子实际上是 $G$ 的内型的那些等晶体）相联系的群（扩张纯内型）一致。对一般的 $G$（例如 $G=\mathrm{GL}_n$），纯内型只对应于基本等晶体的一个子集。

我们在 [BCHN] 中提出：循环形变在所有等晶体与纯内型之间内插，即 $\mathbb L^u_{\underline q}$ 上周期循环层的范畴与 $\check G$ 的纯内型的所有幂幺表示的范畴等同。

\begin{conj}{3.6}\label{conj:a5-02-1}
存在全嵌入
\[
\bigoplus_\eta D_{\mathrm{f.g.}}(\check G_\eta)^u\otimes_k k[u]\subset\Coh\bigl(\mathbb L^u_{\underline q}\bigr)^{S^1},
\]
从 $G$ 的纯内型 $G_\eta$ 的幂幺表示范畴之和的平凡 $u$-形变，到分次幂幺 Langlands 参数上的循环层。此外该嵌入在把 $u$ 可逆化之后成为等价：
\[
\bigoplus_\eta D(\check G_\eta)^u\otimes_k k[u,u^{-1}]\simeq\QCoh^{!}\bigl(\mathbb L^u_{\underline q}\bigr)^{\mathrm{Tate}}.
\]
\end{conj}

换言之，在从所有凝聚层到周期循环层的形变下，我们预期只有纯内型的表示存活，而且形变在这个子范畴上是平凡的。

\begin{rem}{3.7}\label{rem:a5-02-1}
上述猜想中 $S^1$-等变层的嵌入并不期望是等价；即使在 $G=T$ 是环面的情形（例 A.13）也能看出这一点。也就是说，$S^1$-等变范畴有挠，它在把 $u\in C^\bullet(BS^1;k)$ 可逆化时被杀掉。
\end{rem}

\begin{prob}{3.8}\label{prob:a5-02-1}
给出 $S^1$-不变层的满子范畴
\[
\QCoh^{!}\bigl(\mathbb L^u_{\underline q}\bigr)^{\omega S^1}\otimes_{k[u]}k\subset\QCoh^{!}\bigl(\mathbb L^u_{\underline q}\bigr)
\]
的自守刻画。
\end{prob}

\subsection{圆作用从何而来？}
人们自然会想提出一种循环机制，把局部 Langlands 对应的凝聚形式与可构造形式在一般情形下联系起来：一方面推广上面在幂幺非阿基米德情形讨论的结果，另一方面与将在阿基米德情形讨论的结果平行。为此，至少需要在 Langlands 参数的叠上（从而在其凝聚层范畴上）有一个圆作用。

在非阿基米德情形，幂幺算术 Langlands 参数的叠 $\mathbb L^u_q$（$F$ 的 Weil–Deligne 群的表示）是靠在幂幺几何 Langlands 参数的叠 $\NN_{\check G}/\check G$（惯性子群的 Weil–Deligne 表示）上取一个非平凡自同构（对应自守侧的 Frobenius）的导出不动点来实现的，而该叠并不带有圆作用。我们能够这样定义圆作用：转而取 $\NN_{\check G}/(\check G\times\GG_m)$ 的导出圈空间（即恒等映射的导出不动点）的 $S^1$-不变子叠，而后者确实带有这样的作用。在范畴的层次上，恒等算子的迹（Hochschild 同调）带有圆作用，而 Frobenius 的迹则不带。

\ysrc{444}
让我们重新审视幂幺 Langlands 参数叠上圆作用的来源——我们最初是从方程的具体形状认出它的。模空间 $\mathbb L^u_q$ 是「关于 $q$ 的多项式」，意思是它是 $\GG_m$ 上一个自然代数族的纤维，纤维取在 $q\in\GG_m$ 处。当我们以 $\GG_m$-等变的方式工作时，该族的总空间变成一个圈空间——也就是说，我们不是在圈空间上取 $\mathbb L^u_q$ 本身，而是在圈空间上设置 $S^1$-不变的方程「Frobenius 以 $q$ 作用」，从而恢复 $\mathbb L^u_q$ 模去 $\GG_m$ 的商 $\mathbb L^u_{\underline q}$。从实现仿射 Hecke 代数表示的角度看，这个分次版本同样够用，因为凝聚 Springer 层及其所有自同态都是 $\GG_m$-不变的。

\subsubsection{分次幂幺 Langlands 参数}
在幂幺 Langlands 参数的设定中，表达分次提升之来源的另一种方式是通过范畴迹的理论，如 [BNP17, Zh18, Zhu25]。具体地，$\mathbb L^u_q$ 上 ind-凝聚层的整个范畴由 [BNP17] 作为 Frobenius 在仿射 Hecke 范畴（Steinberg 叠上凝聚层的幺半范畴）的谱形式上的作用的范畴迹而出现。取 $K=\overline\F_q((t))$，Bezrukavnikov 的驯顺分歧局部几何 Langlands 对应 [Bez06, Bez16] 给出幺半等价
\[
\DD(I\backslash G_F/I)\simeq\QCoh^{!}\bigl(\wt\NN_{\check G}\times_{\check{\frakg}}\wt\NN_{\check G}\bigr),
\]
它在自守 Hecke 范畴与谱 Hecke 范畴之间把几何 Frobenius 的拉回与乘以 $q$ 的推前缠绕起来，因此两个幺半范畴上 Frobenius 的迹也被等同。在自守侧，Zhu [Zhu25] 把这个迹与 $G(K)$ 的纯内型的幂幺表示（以及在更一般的、与等晶体相联系的群之中的那些）联系起来。

然而 Bezrukavnikov 的等价有一个预期的分次提升（[HL22] 中宣布）。在谱侧，这只需把 $\GG_m$-等变性纳入 Steinberg 叠，于是 Frobenius 的迹变成 $\QCoh^{!}(\mathbb L^u_{\underline q})$。在自守侧，人们得到「混合」仿射 Hecke 范畴 [BY13]，其中我们把反常层换成仿射旗簇上的「Tate」Weil 层，或者更概念地说，换成由 Ho 与 Li [HL22] 在极大一般性下定义的分次 $\ell$-进层范畴。

于是可以把分次仿射 Hecke 范畴上恒等作用的迹——分别在自守侧与谱侧计算——看作提供一种机制，把 [Zhu25] 的结果推广为猜想 3.6 的证明。

\subsubsection{幂幺设定之外}
我们在幂幺设定中的经验提示，应当在 Langlands 参数叠上 $\GG_m$-等变层的范畴上寻找 $S^1$-作用。确实，在 $\mathrm{GL}_n$ 的情形，我们通过标准的归约到 Iwahori 块得到了这样的作用。一般地，Weil–Deligne 表示的叠带有自然的 $\GG_m$-作用（重标幂零自同态），并且可以通过归约到幂幺情形在其商上构造 $S^1$-作用。我们把希望总结为下面这份宽泛的问题清单：

\begin{prob}{3.9}\label{prob:a5-02-2}
把凝聚型与可构造型的非阿基米德局部 Langlands 对应如下联系起来：
\begin{itemize}[leftmargin=1.6em,itemsep=0.1em]
\item 为 Weil–Deligne 表示的叠的 $\GG_m$-商上的 $S^1$-作用找出一个自然的来源。
